\begin{frame}[fragile]{한 뼘 더: 등가 특징 두 개면 어느 쪽을 남기는가}
  \small
  $x_1$과 \emph{완전히 같은 정보}를 담는 등가 특징 $x_{2c}=3x_1-4$(완전
  상관)를 추가 + 잡음 $x_3$:
\begin{lstlisting}
x2c = 3 * x1 - 4.0            # x1과 등가(완전 상관) 특징
x3  = np.array([2,9,1,7,3,6], dtype=float)
Xc  = np.column_stack([x1, x2c, x3])
print(Lasso(alpha=2.0).fit(Xc, y).coef_)
# [0. 0.5905 0.]   (순서: [x1, x2c, x3])
\end{lstlisting}
  놀라운 결과: Lasso는 \textbf{진짜 계수 2를 가진 $x_1$을 0으로 만들고},
  등가인 $x_{2c}$를 남김(계수 0.59, 환산 기울기 $\approx1.77$).
  \vskip0.3em
  ``어느 쪽을 남길지''는 \textbf{데이터가 결정하지 않는다} — 등가이면
  손실함수가 그 방향으로 평탄해져 해가 유일하지 않고, 솔버가 어느 꼭짓점에
  먼저 닿느냐에 따라 둘 중 하나가 \emph{임의로} 낙점.
  \vskip0.2em
  {\footnotesize 대조: Ridge는 상관 특징들의 가중치를 \emph{나눠 갖고}, Lasso는
  하나를 \emph{통째로 삭제} — 안정적 해 vs 희소 해, 둘 다 ``합법적''.}
\end{frame}
